% =============================================================
% บทที่ 5 - สรุปผลการดำเนินงาน
% ปรับปรุงให้ตรงกับสถานะโค้ดจริง ณ ปัจจุบัน (เดือนสิงหาคม 2569) แทนตัวเลขร่างเดิมจาก
% Research Me (1).docx ซึ่งบันทึกไว้ตั้งแต่ช่วงที่ระบบยังแยกส่วนกันอยู่ (ยังไม่เชื่อมต่อ AI จริง)
% หมายเหตุ: ตารางที่ \ref{tab:progress} สรุปเป็น "สิ่งที่ทำสำเร็จแล้ว" / "สิ่งที่ยังคงค้าง" ต่อโมดูล
% (ไม่ใช้ตัวเลข % ความคืบหน้า) ตามสถานะโค้ดจริงในโครงงาน
% =============================================================

% -----------------------------------------------------------
\section{สรุปผลการดำเนินโครงงาน}
% -----------------------------------------------------------

การดำเนินโครงงานในระยะนี้ก้าวข้ามขั้นตอนการทดสอบเทคโนโลยีแกนหลักไปแล้ว โดยระบบทั้งเว็บแอปพลิเคชัน, ระบบหลังบ้าน, และ AI Service เชื่อมต่อและทำงานร่วมกันได้จริงแบบครบวงจรแล้ว ทั้งฟังก์ชันแชตบอตที่ตอบกลับแบบทยอยแสดงผล ฟังก์ชันสร้างและตรวจสอบคุณภาพแผนการเดินทางด้วยกลไก LLM-as-a-Judge อัตโนมัติ และฟังก์ชันเสริมด้าน AI อื่น ๆ เช่น การดึงข้อมูลกิจกรรมจากข้อความโพสต์ และการสร้างคำอธิบายสถานที่เฉพาะราย ทั้งนี้ยังมีบางส่วนที่ยังเป็นต้นแบบ โดยเฉพาะระบบยืนยันตัวตนและระบบสแกน QR สะสมพอยท์ ซึ่งยังไม่ได้เชื่อมต่อกับฐานข้อมูลจริง

ความคืบหน้าของแต่ละส่วนแสดงดังตารางที่ \ref{tab:progress}

{\small
\begin{table}[H]
\centering
\caption{ตารางแสดงผลการดำเนินโครงงาน}
\label{tab:progress}
\begin{tabular}{|p{2.6cm}|p{5.5cm}|p{5.5cm}|}
\hline
\textbf{ส่วนการทำงาน} & \textbf{สิ่งที่ทำสำเร็จแล้ว} & \textbf{สิ่งที่ยังคงค้าง} \\
\hline
1. Backend และฐานข้อมูล & ฐานข้อมูล PostgreSQL (Supabase) 5 ตาราง (places, events, knowledge\_base, rewards, qrs) พร้อม REST API แบบ CRUD ครบทุกตาราง & ระบบยืนยันตัวตนจริงสำหรับผู้ดูแลระบบ; ทุกตารางยังเปิดสิทธิ์อ่าน/เขียนแบบสาธารณะไว้ก่อน \\ \hline
2. User Portal (ผู้ใช้งานทั่วไป) & หน้าหลัก, รายละเอียดสถานที่/กิจกรรม, ค้นหา, แชตบอต, สร้าง-ปรับแผนการเดินทาง, หน้าพอยท์สะสม ครบทุกหน้าจอ เชื่อมต่อ API จริง & สแกน QR/สะสม-แลกพอยท์/เข้าสู่ระบบ ยังเป็น Mock ฝั่ง Frontend ทั้งหมด \\ \hline
3. Admin Portal (ผู้ดูแลระบบ) & จัดการข้อมูลสถานที่ กิจกรรม ฐานความรู้ QR/ของรางวัล แบบ CRUD เชื่อมฐานข้อมูลจริงทั้งหมดแล้ว พร้อมเครื่องมือช่วยดึงข้อมูลกิจกรรมจากโพสต์ด้วย AI & หน้าล็อกอินยังเป็น Mock ไม่มีตรวจสอบสิทธิ์จริง \\ \hline
4. AI \& LLM Integration & RAG (Hybrid Search ผ่าน pgvector), แชตบอต Streaming, วางแผนทริปด้วย LLM พร้อม LLM-as-a-Judge, ดึงข้อมูลกิจกรรมจากโพสต์ด้วย AI เชื่อมเข้าระบบจริงครบแล้ว & การสร้างคำอธิบายสถานที่ด้วย AI ยังเป็นสคริปต์แยกที่ต้องรันเองนอกระบบ ไม่มีปุ่มเรียกใน Admin Portal; กลไก LLM-as-a-Judge ยังไม่แม่นยำเพียงพอ ต้องทดสอบและประเมินผลเพิ่มเติม; ยังไม่ได้ทำ UAT \\ \hline
5. Data Retrieval (ข้อมูลภายนอก) & Google Places API (ข้อมูลตั้งต้น) และ Serper API (ค้นข้อมูลเสริมสร้างคำอธิบายสถานที่) & สำเร็จครบตามที่ออกแบบไว้ \\ \hline
6. Testing \& Deployment & Deploy บริการ AI ขึ้น Docker บน Render แล้ว; มีชุดทดสอบอัตโนมัติสำหรับ Judge และ Trip Planner (Unit Test แบบ Mock คำตอบ LLM) และชุด Golden Set ที่เรียกโมเดลจริงสำหรับฟังก์ชันแชตบอตโดยเฉพาะ & ยังไม่มีชุดทดสอบที่เรียกโมเดลจริงสำหรับ Trip Planner/Judge (มีเฉพาะแบบ Mock); ยังไม่ได้ทำ UAT กับผู้ใช้งานจริง \\
\hline
\end{tabular}
\end{table}
\par}

% -----------------------------------------------------------
\needspace{8\baselineskip}
\section{ข้อจำกัดของระบบ}
% -----------------------------------------------------------

\subsection{ข้อจำกัดด้านการทำงาน}

แม้ระบบหลักจะเชื่อมต่อกันแบบครบวงจรแล้ว แต่ยังพบข้อจำกัดสำคัญในส่วนที่ยังไม่สมบูรณ์ ดังนี้

\begin{enumerate}[label=(\arabic*),leftmargin=0pt,itemindent=\dimexpr\parindent+\parenlabelwidth+1em\relax,align=right,labelwidth=\parenlabelwidth,labelsep=1em]
\needspace{5\baselineskip}
\item ยังไม่มีระบบยืนยันตัวตนจริง ทั้งฝั่งผู้ดูแลระบบและฝั่งนักท่องเที่ยว หน้าล็อกอินที่ใช้สาธิตอยู่เป็นเพียงการจำลองสถานะฝั่ง Frontend เท่านั้น ยังไม่มีการตรวจสอบรหัสผ่านหรือสิทธิ์การเข้าถึงจากฝั่ง Backend และตารางฐานข้อมูลทุกตารางยังเปิดสิทธิ์อ่าน/เขียนแบบสาธารณะไว้ก่อน
\item ระบบสแกน QR และสะสม/แลกพอยท์ ยังเป็นต้นแบบ การผูก QR กับสถานที่และรายการของรางวัลจัดการผ่านฐานข้อมูลจริงได้แล้วในฝั่งผู้ดูแลระบบ แต่ขั้นตอนฝั่งผู้ใช้งานยังเป็นต้นแบบทั้งหมด: ปุ่ม ``สแกน QR'' จำลองผลลัพธ์ด้วยการสุ่มเลือกสถานที่ที่มี QR จากรายการที่โหลดไว้แล้ว ยังไม่ได้เปิดกล้องอ่านรหัสจริง ยังไม่มีการตรวจสอบการเข้าสู่ระบบหรือจำกัดความถี่การสแกนใด ๆ และพอยท์สะสมยังเก็บอยู่ในหน่วยความจำของเบราว์เซอร์ ไม่ได้ผูกกับบัญชีผู้ใช้หรือบันทึกลงฐานข้อมูล ส่วนขั้นตอนแลกของรางวัล แม้รายการของรางวัลและราคาพอยท์จะดึงมาจากฐานข้อมูลจริงแล้ว แต่การตัดพอยท์เมื่อแลกยังเป็นปุ่มจำลองอยู่ที่หน้าพอยท์สะสมฝั่งผู้ใช้งาน ยังไม่ได้ย้ายไปยัง Admin Portal ตามที่ออกแบบไว้ และยังไม่มีการบันทึกประวัติการแลกหรือสถานะการรับของรางวัลลงฐานข้อมูล
\item แผนการเดินทางยังไม่ถูกบันทึกถาวร ระบบยังไม่มีตาราง Itinerary สำหรับเก็บประวัติแผนการเดินทางที่ AI สร้างให้ ผู้ใช้งานจึงยังไม่สามารถย้อนดูแผนเก่า หรือให้ผู้ดูแลระบบตรวจสอบแผนที่เคยสร้างได้ การปรับแผน ก็ดำเนินการอยู่ในหน่วยความจำของเบราว์เซอร์เช่นกัน โดยเลือกสถานที่สำรองจากรายการที่ดึงมาแล้วเท่านั้น ไม่ได้เรียกให้ LLM สร้างแผนใหม่ซ้ำ
\end{enumerate}

\subsection{ข้อจำกัดด้านฮาร์ดแวร์และซอฟต์แวร์}

ระบบพึ่งพาการประมวลผลผ่านโมเดลภาษาขนาดใหญ่ตระกูล Gemini ที่เรียกผ่าน API Gateway ของมหาวิทยาลัยขอนแก่น ซึ่งมีการกำหนดเกณฑ์การใช้งาน ทั้งในแง่ของจำนวนครั้งต่อนาทีและจำนวน Token ต่อวัน ทำให้อาจเกิดการหยุดชะงักของบริการเมื่อโควต้าหมดลงหากมีผู้ใช้งานพร้อมกันจำนวนมาก นอกจากนี้ ค่าพารามิเตอร์บางส่วนของกลไก LLM-as-a-Judge ถูกปรับตั้งจากผลทดสอบกับโมเดลชุดแรกที่ใช้งานเท่านั้น หากเปลี่ยนไปใช้โมเดลรุ่นอื่นในอนาคต จำเป็นต้องทดสอบและปรับค่าดังกล่าวใหม่ก่อนนำไปใช้งานจริง

% -----------------------------------------------------------
\needspace{6\baselineskip}
\section{ปัญหาอุปสรรค และแนวทางแก้ไข}
% -----------------------------------------------------------

ในระหว่างการพัฒนาและทดสอบระบบ คณะผู้จัดทำพบปัญหาทางเทคนิคหลายประการ โดยเฉพาะในส่วนของการรวบรวมข้อมูลเชิงลึกจากแหล่งข้อมูลภายนอกและการควบคุมการใช้งานโควต้าของ LLM ซึ่งได้มีการวิเคราะห์สาเหตุและดำเนินการแก้ไขบางส่วนไปแล้ว ดังแสดงในตารางที่ \ref{tab:problems}

\begin{table}[H]
\centering
\caption{ตารางแสดงการสรุปปัญหา อุปสรรค และแนวทางแก้ไข}
\label{tab:problems}
\begin{tabular}{|p{3.8cm}|p{4.2cm}|p{5.5cm}|}
\hline
\textbf{ปัญหาและอุปสรรค} & \textbf{สาเหตุหลัก} & \textbf{แนวทางแก้ไข} \\
\hline
ไม่สามารถดึงข้อมูล Popular Times จาก Google Maps ได้อย่างสม่ำเสมอ & การเขียนสคริปต์ดึงข้อมูล (Web Scraping) ถูกบล็อกโดยระบบป้องกันการดึงข้อมูลอัตโนมัติ (Anti-Scraping / IP Blocking) ของ Google & ยังเป็นข้อจำกัดที่ไม่ได้รับการแก้ไข ปัจจุบันระบบใช้ข้อมูลรีวิว เวลาเปิด-ปิด และสถานะการเปิดให้บริการ (business\_status) จาก Google Places API เป็นตัวแปรทดแทนไปก่อน ยังไม่ได้ทดลองใช้ Third-Party API สำหรับข้อมูลความหนาแน่นโดยเฉพาะ \\ \hline
โควต้าของ LLM อาจหมดลงได้หากมีผู้ใช้งานพร้อมกันจำนวนมาก & ระบบพึ่งพาการประมวลผลผ่าน Gemini API ที่เรียกผ่าน KKU Gateway ซึ่งมีการกำหนดเกณฑ์การใช้งาน (Rate Limits) & 1. แยกค่าโมเดลตามลักษณะงาน (แชตบอต/Event Extraction ใช้โมเดลขนาดเล็กกว่า Trip Planner Generator/Judge) เพื่อกระจายภาระและควบคุมต้นทุนต่อคำขอ \newline 2. จำกัดจำนวนรอบการสร้างแผนซ้ำของ LLM-as-a-Judge ไว้ไม่เกินค่าที่กำหนด (MAX\_JUDGE\_ATTEMPTS) เพื่อไม่ให้คำขอเดียวใช้โควต้าเกินความจำเป็น \newline 3. ตั้งค่าให้ขั้นตอนตรวจสอบคุณภาพ (Judge) ทำงานแบบ Fail-Open หากเรียก API ไม่สำเร็จ เพื่อไม่ให้ปัญหาโควต้าที่ขั้นตอนตรวจสอบไปปิดกั้นผลลัพธ์หลักที่สร้างเสร็จแล้ว \newline (แนวทาง Local LLM/Semantic Caching ที่เคยเสนอไว้ในรายงานฉบับก่อนหน้ายังไม่ได้นำมาใช้จริง) \\
\hline
\end{tabular}
\end{table}

% -----------------------------------------------------------
\needspace{6\baselineskip}
\section{ข้อเสนอแนะ ในการพัฒนาต่อไป}
% -----------------------------------------------------------

\needspace{6\baselineskip}
\subsection{การปรับปรุงฟังก์ชันการทำงาน}

เพื่อให้ระบบมีความสมบูรณ์และพร้อมใช้งานจริง ควรเร่งดำเนินการพัฒนาระบบยืนยันตัวตนที่แท้จริงทั้งฝั่งผู้ดูแลระบบและฝั่งนักท่องเที่ยว รวมถึงเชื่อมต่อระบบสแกน QR สะสม/แลกพอยท์เข้ากับฐานข้อมูลจริงแทนการจำลองผลในเบราว์เซอร์ และเพิ่มตารางสำหรับบันทึกประวัติแผนการเดินทางที่สร้างไว้ นอกจากนี้ ควรศึกษาแนวทางการเข้าถึงข้อมูลความหนาแน่นของผู้คนจากแหล่งข้อมูลเปิดอื่นที่ยังไม่เคยทดลอง เพื่อเพิ่มประสิทธิภาพในการแนะนำแผนการเดินทางให้สอดคล้องกับสถานการณ์จริงยิ่งขึ้นและเมื่อฟังก์ชันหลักครบถ้วนสมบูรณ์ ควรนำระบบไปทดสอบกับผู้ใช้งานจริงด้วย UAT เพื่อประเมินประสิทธิภาพและความพึงพอใจของผู้ใช้งาน รวมถึงทดสอบซ้ำและปรับค่าพารามิเตอร์ของกลไก LLM-as-a-Judge หากมีการเปลี่ยนโมเดลภาษาที่ใช้งาน

\subsection{แนวทางการวิจัยต่อยอด}

จากผลการดำเนินงานและปัญหาอุปสรรคที่พบ คณะผู้จัดทำได้กำหนดแผนการดำเนินงานในช่วงถัดไป เพื่อให้โครงงานสามารถพัฒนาเสร็จสมบูรณ์และนำไปใช้งานได้จริง โดยมีรายละเอียดดังแสดงในตารางที่ \ref{tab:nextplan}

\begin{table}[H]
\centering
\caption{แผนการดำเนินงานในระยะต่อไป}
\label{tab:nextplan}
\begin{tabular}{|p{2.4cm}|p{3.6cm}|p{7cm}|}
\hline
\textbf{ระยะเวลา} & \textbf{แผนการดำเนินงาน} & \textbf{รายละเอียด} \\
\hline
เดือนสิงหาคม 2569 & พัฒนาโค้ดให้เสร็จสมบูรณ์ตาม Use Case & พัฒนาระบบยืนยันตัวตนจริงสำหรับผู้ดูแลระบบและนักท่องเที่ยว (UC-11); เชื่อมระบบสแกน QR และระบบสะสม/แลกพอยท์เข้ากับฐานข้อมูลจริงแทนการจำลองผลในเบราว์เซอร์ (UC-03, UC-08); เพิ่มตารางบันทึกประวัติแผนการเดินทาง (Itinerary); เพิ่มการตรวจสอบความสัมพันธ์ระหว่างตารางก่อนลบข้อมูลสถานที่ที่มี QR ผูกอยู่ ให้ครบตามขอบเขตของทุก Use Case ที่ออกแบบไว้ \\ \hline
ต้นเดือนกันยายน 2569 & Deploy ระบบขึ้นใช้งานจริง & นำระบบที่พัฒนาเสร็จสมบูรณ์แล้วขึ้น Deploy บนสภาพแวดล้อมจริง พร้อมปรับสิทธิ์การเข้าถึงฐานข้อมูล (Row Level Security) ให้จำกัดเฉพาะผู้ใช้ที่ผ่านการยืนยันตัวตนแล้วก่อนเริ่มขั้นตอนการทดสอบ \\ \hline
ปลายเดือนกันยายน 2569 & ทดสอบระบบและเปรียบเทียบโมเดล LLM & ทดสอบระบบที่ Deploy แล้วกับผู้ใช้งานจริงด้วย UAT เพื่อประเมินประสิทธิภาพและความพึงพอใจของผู้ใช้งาน ควบคู่กับการเปรียบเทียบโมเดลภาษาขนาดใหญ่ (LLM) หลายรุ่น เพื่อประเมินความแม่นยำ ความเร็วในการตอบสนอง และต้นทุนการใช้งาน ก่อนคัดเลือกโมเดลที่เหมาะสมที่สุดสำหรับแต่ละฟังก์ชัน (แชตบอต, Trip Planner, Judge) รวมถึงทดสอบซ้ำและปรับค่าพารามิเตอร์ของกลไก LLM-as-a-Judge ให้เหมาะกับโมเดลที่เลือกใช้จริง \\ \hline
เดือนตุลาคม 2569 & วิเคราะห์ผลและจัดทำเอกสารฉบับสมบูรณ์ & วิเคราะห์ผลการทดสอบระบบ และข้อเสนอแนะที่ได้รับจากผู้ใช้งาน สรุปผลการเปรียบเทียบโมเดล LLM ที่ทดสอบ จัดทำรายงานโครงงานฉบับสมบูรณ์และเตรียมนำเสนอ \\
\hline
\end{tabular}
\end{table}
